\chapter{Tín hiệu báo động}
% full version: chapters/14_pain.tex

\pencilfig{14}{\pencilart{14}}{Cơn đau là chuông báo cháy, không phải nhiệt kế.}

Mọi giác quan khác báo thế giới ra sao. Cơn đau thì không. Nó báo rằng có gì đó đang đe dọa cơ thể, và độ to của nó phần lớn do chính cỗ máy đặt ra chứ không phải do kích thích. Cơn đau là chuông báo cháy, không phải nhiệt kế.

Cảm biến đau chỉ kích hoạt với cái mạnh: với nhiệt trên bốn mươi mấy độ, với áp lực mạnh, với hóa chất ăn mòn. Chúng quen dần yếu hơn nhiều so với cảm biến xúc giác --- chuông báo động không được im khi mối đe dọa còn đó.

\section*{Hai dây dẫn}

Hãy nhớ lại lúc bạn bị va ngón tay. Trước tiên là cơn đau nhói, chính xác: ở đâu và khi nào. Một khoảnh khắc sau là cơn đau âm ỉ, lan tỏa, kéo dài: tệ đến mức nào. Đó là hai dây dẫn khác nhau. Một dây nhanh, tín hiệu của nó bay tới trong vài chục mili giây. Dây kia chậm, tín hiệu của nó đi mất cỡ một giây.

\section*{Cánh cổng}

Vì sao khi bị va, ta lại xoa chỗ đau? Trong tủy sống, tín hiệu đau đi qua một ``cánh cổng''. Cú chạm đánh thức một nơ-ron ức chế, và nơ-ron này khép cổng lại. Vì thế xoa chỗ va quả thật có ích --- nhưng chỉ với cơn đau vừa phải; cơn đau mạnh vẫn đi qua. Còn việc cơn đau mạnh khi đó tự tắt nơ-ron ức chế là một phần của lý thuyết cổng ban đầu, đến nay chưa có bằng chứng trực tiếp.

\pencilfig{126}{%
% gate: the pain wire drives the relay neuron; touch wakes an inhibitory neuron that closes the gate
\draw[pen=pcAmber] (0,1.35) -- (3.05,1.0); \draw[arr=pcAmber] (2.8,1.03) -- (3.08,1.0);
\node[lbl, text=pcAmber!70!black, anchor=east] at (-0.05,1.35) {đau};
\draw[pen=pcBlue] (0,-0.15) -- (1.75,-0.15); \draw[arr=pcBlue] (1.5,-0.15) -- (1.78,-0.15);
\node[lbl, text=pcBlue, anchor=east] at (-0.05,-0.15) {cú chạm};
\draw[dotpen=pcRed, wash=pcRed] (2.05,-0.15) circle (0.26);
\draw[pcRed, line width=0.7pt, -{Bar[width=6pt]}] (2.25,0.05) -- (3.2,0.66);
\node[lbl, text=pcRed, anchor=north] at (2.05,-0.45) {phanh};
\draw[dotpen=pcGray, wash=pcGray] (3.4,0.9) circle (0.34);
\node[lbl, anchor=south] at (3.4,1.3) {cổng};
\draw[pen] (3.75,0.9) -- (5.6,0.9); \draw[arr] (5.35,0.9) -- (5.65,0.9);
\node[lbl, anchor=west] at (5.7,0.9) {lên não};
}{Cánh cổng trong tủy sống. Tín hiệu đau đi lên não qua nơ-ron cổng, còn cú chạm đánh thức một nơ-ron ức chế khép cổng lại. Vì thế xoa chỗ va có ích, nhưng chỉ với cơn đau vừa phải.}

\section*{Núm âm lượng từ trên xuống}

Bộ não tự biết làm cơn đau to lên hay nhỏ đi. Từ trên xuống, từ não bộ tới cánh cổng, có các kênh radio --- \nt{SERO}, \nt{NORA} và những chất đặc biệt giống thuốc giảm đau do chính bộ não tiết ra. Trong tình huống nguy hiểm, cơn đau có thể gần như biến mất: người đang chạy trốn hiểm nguy không có thời gian để khập khiễng. Sự chờ đợi được giảm đau cũng làm dịu cơn đau, một phần qua chính những chất đó.

\section*{Khi chuông báo động bị kẹt}

Cơn đau lặp lại làm tăng độ nhạy: cùng một cú va bắt đầu đau hơn. Đó cũng là một kiểu học --- một vòng lặp tự vặn to mình lên trong lúc vết thương lành. Nhưng một vòng lặp mạnh có một mối nguy: nó có thể ``kẹt'' và tiếp tục reo khi vết thương đã lành từ lâu --- như nhóm nơ-ron ở chương hai tự duy trì chính mình.

\section*{Thầy giáo và ngắt}

Cơn đau còn là một thầy giáo: với tín hiệu dopamine ``tốt hơn hay tệ hơn mong đợi'', nó hoạt động như một phần thưởng âm. Và nó ngắt việc đang làm, như một đợt bùng noradrenaline: bỏ hết đi, lo mối nguy trước. Còn phản ứng nhanh nhất --- rụt tay khỏi vật nóng --- khép kín ngay trong tủy sống, nhờ chính cặp cơ đối kháng ở chương trước, trước cả khi bạn kịp cảm thấy đau.

\fullversion{14}
